 \documentclass[12pt,letterpaper]{article}

% Package imports
\usepackage{latexsym}
\usepackage{xcolor}
\usepackage{float}
\usepackage{ragged2e}
\usepackage[empty]{fullpage}
\usepackage{wrapfig}
\usepackage{lipsum}
\usepackage{tabularx}
\usepackage{titlesec}
\usepackage{geometry}
\usepackage{marvosym}
\usepackage{verbatim}
\usepackage{enumitem}
\usepackage{fancyhdr}
\usepackage{multicol}
\usepackage{graphicx}
\usepackage{cfr-lm}
\usepackage[T1]{fontenc}
\usepackage{fontawesome5}

\definecolor{darkblue}{RGB}{0,0,139}

% Page layout
\setlength{\multicolsep}{0pt} 
\pagestyle{fancy}
\fancyhf{} % clear all header and footer fields
\fancyfoot{}
\renewcommand{\headrulewidth}{0pt}
\renewcommand{\footrulewidth}{0pt}
\geometry{margin=0.75in}
\setlength{\footskip}{5pt} % Addressing fancyhdr warning

% Hyperlink setup (moved after fancyhdr to address warning)
\usepackage[hidelinks]{hyperref}
\hypersetup{
    colorlinks=true,
    linkcolor=darkblue,
    filecolor=darkblue,
    urlcolor=darkblue,
}

% Custom box settings
\usepackage[most]{tcolorbox}
\tcbset{
    frame code={},
    center title,
    left=0pt,
    right=0pt,
    top=0pt,
    bottom=0pt,
    colback=gray!20,
    colframe=white,
    width=\dimexpr\textwidth\relax,
    enlarge left by=-2mm,
    boxsep=4pt,
    arc=0pt,outer arc=0pt,
}

% URL style
\urlstyle{same}

% Text alignment
\raggedright
\setlength{\tabcolsep}{0in}

% Section formatting
\titleformat{\section}{
  \vspace{-4pt}\scshape\raggedright\large
}{}{0em}{}[\color{black}\titlerule \vspace{-7pt}]

% Custom commands
\newcommand{\resumeItem}[2]{
  \item{
    \textbf{#1}{\hspace{0.5mm}#2 \vspace{-0.5mm}}
  }
}

\newcommand{\resumePOR}[3]{
\vspace{0.5mm}\item
    \begin{tabular*}{0.97\linewidth}[t]{l@{\extracolsep{\fill}}r}
        \textbf{#1}\hspace{0.3mm}#2 & \textit{\small{#3}} 
    \end{tabular*}
    \vspace{-2mm}
}

\newcommand{\resumeSubheading}[5]{
\vspace{0.5mm}\item
    \begin{tabular*}{0.98\linewidth}[t]{l@{\extracolsep{\fill}}r}
        \textbf{#1} & \textit{\footnotesize{#5}} \\
        {\small{\textit{#3}}\hspace{1em}\small{#4}} & \footnotesize{#2}\\
    \end{tabular*}
    \vspace{-2.4mm}
}

\newcommand{\resumeProject}[4]{
\vspace{0.5mm}\item
    \begin{tabular*}{0.98\linewidth}[t]{l@{\extracolsep{\fill}}r}
        \textbf{#1} & \textit{\footnotesize{#3}} \\
        \footnotesize{\textit{#2}} & \small{#4}
    \end{tabular*}
    \vspace{-2.4mm}
}

\newcommand{\resumeOneLine}[2]{
\vspace{0.5mm}\item
    \begin{tabular*}{0.98\linewidth}[t]{l@{\extracolsep{\fill}}r}
        \textbf{#1} & \textit{\footnotesize{#2}}
    \end{tabular*}
    \vspace{-1.2mm}
}

\newcommand{\resumeSubItem}[2]{\resumeItem{#1}{#2}\vspace{-4pt}}

\renewcommand{\labelitemi}{$\vcenter{\hbox{\tiny$\bullet$}}$}
\renewcommand{\labelitemii}{$\vcenter{\hbox{\tiny$\circ$}}$}

\newcommand{\resumeSubHeadingListStart}{\begin{itemize}[leftmargin=*,labelsep=3mm]}
\newcommand{\resumeHeadingSkillStart}{\begin{itemize}[leftmargin=*,itemsep=1.7mm, rightmargin=2ex]}
\newcommand{\resumeItemListStart}{\begin{itemize}[leftmargin=*,labelsep=3mm,itemsep=0.5mm]}

\newcommand{\resumeSubHeadingListEnd}{\end{itemize}\vspace{-1mm}}
\newcommand{\resumeHeadingSkillEnd}{\end{itemize}\vspace{-1mm}}
\newcommand{\resumeItemListEnd}{\end{itemize}\vspace{-1mm}}
\newcommand{\cvsection}[1]{%
\vspace{2mm}
\begin{tcolorbox}
    \textbf{\large #1}
\end{tcolorbox}
    \vspace{-4mm}
}

\newcolumntype{L}{>{\raggedright\arraybackslash}X}%
\newcolumntype{R}{>{\raggedleft\arraybackslash}X}%
\newcolumntype{C}{>{\centering\arraybackslash}X}%

% Commands for icon sizing and positioning
\newcommand{\socialicon}[1]{\raisebox{-0.05em}{\resizebox{!}{1em}{#1}}}
\newcommand{\ieeeicon}[1]{\raisebox{-0.3em}{\resizebox{!}{1.3em}{#1}}}

% Font options
\newcommand{\headerfonti}{\fontfamily{phv}\selectfont} % Helvetica-like (similar to Arial/Calibri)
\newcommand{\headerfontii}{\fontfamily{ptm}\selectfont} % Times-like (similar to Times New Roman)
\newcommand{\headerfontiii}{\fontfamily{ppl}\selectfont} % Palatino (elegant serif)
\newcommand{\headerfontiv}{\fontfamily{pbk}\selectfont} % Bookman (readable serif)
\newcommand{\headerfontv}{\fontfamily{pag}\selectfont} % Avant Garde-like (similar to Trebuchet MS)
\newcommand{\headerfontvi}{\fontfamily{cmss}\selectfont} % Computer Modern Sans Serif
\newcommand{\headerfontvii}{\fontfamily{qhv}\selectfont} % Quasi-Helvetica (another Arial/Calibri alternative)
\newcommand{\headerfontviii}{\fontfamily{qpl}\selectfont} % Quasi-Palatino (another elegant serif option)
\newcommand{\headerfontix}{\fontfamily{qtm}\selectfont} % Quasi-Times (another Times New Roman alternative)
\newcommand{\headerfontx}{\fontfamily{bch}\selectfont} % Charter (clean serif font)

\begin{document}
\headerfontii

% Header
\begin{center}
    {\Huge\textbf{GUANYU XU}}
\end{center}
\vspace{-6mm}

\begin{center}
    \small{
    +1-734-330-1367 | \href{mailto:xuguanyu@umich.edu}{xuguanyu@umich.edu} | 
    \href{https://xuguaaaanyu.github.io/}{Personal Homepage}
    }
\end{center}
\vspace{-6mm}

%\begin{center}
%    \small{
%    \socialicon{\faLinkedin} \href{https://www.linkedin.com/in/your-profile/}{your-profile} | 
%    \socialicon{\faGithub} \href{https://github.com/your-username}{your-username} | 
%    \ieeeicon{\includegraphics[height=1.3em]{ieee_collabratec_icon.png}} \href{https://ieee-collabratec.ieee.org/app/p/your-profile}{your-profile} |
%    \socialicon{\faTwitter} \href{https://twitter.com/your-handle}{your-handle}
%    }
%\end{center}
%\vspace{-6mm}
\begin{center}
    \small{Ann Arbor, MI - 48105, United States}
\end{center}

\vspace{-4mm}

\section{\textbf{Education}}
\vspace{-0.4mm}
\resumeSubHeadingListStart

\resumeSubheading
{University of Michigan - Ann Arbor}{Ann Arbor, US}
{B.S.E. in Computer Engineering (Expected)}
{\textbf{Overall GPA}: 3.97/4.0, \textbf{Major GPA}: 4.0/4.0}
{Aug 2024 - May 2026}

\resumeSubheading
{Shanghai Jiao Tong University}{Shanghai, China}
{B.E. in Mechanical Engineering (Expected)}{\textbf{Overall GPA}: 3.63/4.0, \textbf{Major GPA}: 3.68/4.0}{Sep 2022 - Aug 2026}
\resumeSubHeadingListEnd

\section{\textbf{Research Interest}}
\small{
My research interest primarily focuses on \textbf{embodied intelligence and robotics}, with the long-term goal of enabling robots to perform diverse tasks in unstructured human environments. I am particularly interested in unifying robotic foundation models with simulation-based RL to enable agile, generalizable robotic behaviors in the real world. My prior work is grounded in \textbf{full-stack robotic systems}, encompassing expertise in mechanical design and smart manufacturing, embedded and real-time systems, and closed-loop sensorimotor control.
\section{\textbf{Publication \& Patents}}
\small{
\begin{enumerate}[leftmargin=*, labelsep=0.5em, align=left, widest={[\textbf{S.1}]}, itemindent=0em, label={\textbf{[\arabic*]}]}]
%\item[\textbf{[C.1]}] Your Name, et al. (Year). \href{https://doi.org/XX.XXXX/XXXXXXX.XXXX.XXXXXXX}{\textbf{Title of Conference Paper}}. In \textit{Name of Conference Proceedings}, pp. XX-XX. Publisher. Date, Location. DOI: XX.XXXX/XXXXXXX.XXXX.XXXXXXX

%\item[\textbf{[S.1]}] Your Name, et al. (Year). \textbf{Title of Submitted Paper}. Manuscript submitted for publication in \textit{Journal Name}.

\item[\textbf{[1]}] \textbf{\underline{Guanyu Xu}}, Jiaqi Wang, Dezhong Tong, and Xiaonan Huang, "Highly Deformable Proprioceptive Membrane for Real-Time 3D Shape Reconstruction", \textit{ArXiv.org. \href{https://arxiv.org/abs/2601.13574}{https://arxiv.org/abs/2601.13574}}, 2026. 

\item[\textbf{[2]}] \textbf{\underline{Guanyu Xu}} and Longquan Liu, "A Variable Radius Wheel", \textit{National Intellectual Property Office, Patent No. \href{http://epub.cnipa.gov.cn/cred/CN221698395U}{
CN221698395U}}, 2024.
%\item[\textbf{[J.1]}] Author 1, Your Name, Author 3, et al. (Year). \href{https://doi.org/XX.XXXX/XXXXX.XXXX.XXXXXXX}{\textbf{Title of Journal Article}}. \textit{Journal Name}, Vol. XX, Issue X, pp. XXX-XXX. DOI: XX.XXXX/XXXXX.XXXX.XXXXXXX
\end{enumerate}
}

\section{\textbf{Research Experience}}
\resumeSubHeadingListStart
% Optical Sensor Project
\resumeProject
    {Proprioceptive Membrane for 3D Shape Reconstruction (Project lead)}
    {\href{https://soft.robotics.umich.edu/}{\underline{Hybrid Dynamic Robotics Lab}}, University of Michigan}
    {Jun. 2025 - Present}
    {Instructor: Prof. Xiaonan Huang}
\vspace{2mm}
\resumeItemListStart
    \item \textbf{Spearheaded} the design and fabrication of a highly deformable proprioceptive membrane utilizing an \textbf{optical waveguide structure} to achieve real-time \textbf{3D shape reconstruction}.
    \item Architected and trained a \textbf{PointNet-based autoencoder model} in \textbf{PyTorch} to decode optical signals into 3D point clouds.
    \item Developed an automated data collection pipeline with a \textbf{depth camera} to capture high-accuracy ground-truth datasets.
    \item Programmed \textbf{STM32} microcontroller to scan LEDs and sample photodiodes, and designed a custom stretchable PCB for the multilayer optical waveguide structure. 
    \item Achieved high-accuracy surface shape reconstruction with \textbf{an average Chamfer distance of 1.3 mm} while maintaining accuracy for \textbf{indentations up to 25 mm} and \textbf{up to 75\% stretching strain}.
    \item \textbf{First-authored paper under review at \textit{Advanced Robotics Research}.}
\resumeItemListEnd

% Growing Robot Project
\resumeProject
    {Active Steering Control of Soft Growing Robot (Core contributor)}
    {\href{https://soft.robotics.umich.edu/}{\underline{Hybrid Dynamic Robotics Lab}}, University of Michigan}
    {Oct. 2024 - Mar. 2025}
    {Instructor: Prof. Xiaonan Huang}
\vspace{2mm}
\resumeItemListStart
    \item Integrated the electrostatic clutch-based steering joint into a functional full-scale robot prototype to enable accurate closed-loop omnidirectional steering. 
    \item Formulated a \textbf{geometrical model} to characterize the non-linear relationship between the steering angle and the clutch actuation patterns.
    \item Designed a custom PCB to drive the high-voltage electrostatic clutch control circuits reliably.
\resumeItemListEnd

% Transformable Wheel Project
\resumeProject
    {Lunar Rover with Transformable Wheel (Project lead)}
    {School of Aeronautics and Astronautics, Shanghai Jiao Tong University}
    {Jan. 2023 - Oct. 2023}
    {Instructor: Prof. Longquan Liu}
\vspace{2mm}
\resumeItemListStart
    \item \textbf{Invented and prototyped} a self-adaptive, variable-radius transformable wheel, successfully translating the design into a fully granted \textbf{National Patent}.
    \item Implemented real-time sensing and transformation actuation algorithms on a Raspberry Pi, interfacing with ultrasonic sensors and LiDAR for obstacle-aware terrain adaptation. 
    \item Developed and tuned a PID controller utilizing IMU feedback to maintain vehicle path stabilization during dynamic wheel transformations. 
\resumeItemListEnd
\resumeSubHeadingListEnd

\section{\textbf{Project Experience}}
\resumeSubHeadingListStart
% INSIGHT Project
\resumeProject
  {INSIGHT: Smart Assistive Glass for Visually Impaired (Team Lead)}
  {Course project for \href{https://www.eecs.umich.edu/courses/eecs473/overview.html}{\underline{EECS 473: Advanced Embedded System}} (MDE).}
  {Aug. 2025 - Dec. 2025}
  {Instructor: Prof. Mark Brehob}
\vspace{2mm}
\resumeItemListStart
    \item Designed an \textbf{ESP32-based smart-glasses system} on a \textbf{custom PCB}, integrating microphone arrays, speaker, vibration motors, and OV2640 camera in a compact wearable form factor.
    \item Engineered an end-to-end edge AI pipeline leveraging a \textbf{Jetson Orin Nano} base station, enabling real-time obstacle-aware assistive navigation and interactive scene description. 
    \item Deployed \textbf{YOLO object detection} on the base station to achieve \textbf{up to 30 Hz} obstacle identification with directional vibrotactile feedback for assistive navigation. 
    \item Deployed a lightweight Vision-Language Model (VLM) for strictly \textbf{on-device} scene description, performing real-time inference without relying on cloud APIs. 
    \item Programmed the ESP32 to handle wake-word detection, real-time audio and video streaming, and low-latency text-to-speech (TTS) playback. 
\resumeItemListEnd

% LuminGrid Project
\resumeProject
    {Lumen Grid: Multi-Robot Competitive Parking Game (Team Lead)}
    {Course project for \href{https://www.eecs.umich.edu/courses/eecs373/index.html}{\underline{EECS 373: Introduction to Embedded System Design}}}
    {Feb. 2025 - Apr. 2025}
    {Instructor: Prof. Junyi Zhu}
\vspace{2mm}
\resumeItemListStart
    \item Programmed \textbf{robot control logic and inter-system communication protocols} in \textbf{C++} on an STM32 microcontroller.
    \item Designed an \textbf{IMU-based remote controller} for the \textbf{Zumo robot} with vibration feedback reflecting Zumo’s speed.
    \item Interfaced with a camera for real-time position tracking of all robots based on color codes.
    \item Developed the \textbf{playground control algorithm} that scheduled lighting patterns, tracked robot position, and updated the scoring for the game setting.
\resumeItemListEnd
\resumeSubHeadingListEnd

\section{\textbf{Teaching \& Community Service Experience}}
\vspace{-0.4mm}
\resumeSubHeadingListStart
  \resumeOneLine{Grader, EECS 370: Intro to Computer Organization, University of Michigan}{Sep. 2025 -- Dec. 2025}
  \resumeOneLine{Volunteer, IEEE International Conference on Robotics and Automation, Atlanta}{May. 2025}
  \resumeOneLine{Student Volunteer, Shanghai Sunflower Community Children’s Service Center}{Sep. 2023 -- Dec. 2023}
\resumeSubHeadingListEnd

\section{\textbf{Honors \& Awards}}
\vspace{-0.4mm}
\resumeSubHeadingListStart
  \resumeOneLine{University Honors, University of Michigan}{Dec. 2024 \& May 2025}
  \resumeOneLine{Dean's List, University of Michigan}{Dec. 2024 \& Apr. 2025}
  \resumeOneLine{Undergraduate Excellent Scholarship, Shanghai Jiao Tong University}{Dec 2023}
\resumeSubHeadingListEnd


\section{\textbf{Skills}}
\vspace{-0.4mm}
\resumeHeadingSkillStart
  \resumeItem{Programming: }{C/C++, Python, MATLAB, Bash, CUDA C/C++.}
  \resumeItem{Embedded Systems: }{STM32, ESP32; Raspberry Pi, Jetson Orin Nano; FreeRTOS; I$^2$C/SPI/UART/CAN. }
  \resumeItem{AI/ML: }{PyTorch, YOLO, Moondream-VLM, Sentence Transformer.}
  \resumeItem{Hardware Design: }{Altium Designer, EasyEDA, Solidworks, AutoCAD.}
  \resumeItem{Tools: }{Git, SSH, CMake, Vim, HTML, Markdown, \LaTeX.}
\resumeHeadingSkillEnd
\end{document}